% 《大学物理简明教程》第5版，第11章狭义相对论，填空11.2.1—11.2.6。
% 原书/扫描PDF第289—290页，原题短摘要及推导独立撰写。
\subsection{填空题：基本原理、时空变换与动能}

\sourceproblem{11.2.1}{289}{289}
\problemtext 写出狭义相对论两条基本公设：物理规律的惯性系等价性以及真空光速。
\workings
第一空是\textbf{所有惯性系}；第二空是\textbf{等于}。
相对性原理要求物理规律在所有惯性系具有相同形式；
光速不变原理给所有惯性系中真空光速均为
\[
\boxed{c=299\,792\,458\,\mathrm{m/s}}.
\]
\crosscheck 光速不变不等于光速在任意物质介质中不变。

\sourceproblem{11.2.2}{290}{290}
\problemtext 飞船沿$x$正向匀速前进，船首和船尾各有发光装置。
飞船尾部观察者测量由船首发出的光脉冲传播速度大小，
船首观察者测量由船尾发出脉冲的传播速度大小，分别填空。
\workings
两个接收器都是惯性系中的局部观测者，
真空光速均满足
\[
\boxed{u_{\mathrm{head\to tail}}=c},\qquad
\boxed{u_{\mathrm{tail\to head}}=c}.
\]
\crosscheck 不能使用伽利略式$c\pm v$；只要是在真空中进行局域光速测定，结果与飞船速度无关。

\sourceproblem{11.2.3}{290}{290}
\problemtext $S'$以$0.8c$沿$S$的$x$正向运动，
$S'$系中两事件相隔
$\Delta t'=5\times10^{-7}$ s、$\Delta x'=-120$ m，
求$S$系中空间与时间间隔。
\workings
逆洛伦兹变换
\[
\Delta x=\gamma(\Delta x'+v\Delta t'),
\quad
\Delta t=\gamma\left(\Delta t'+\frac{v\Delta x'}{c^2}\right),
\quad\gamma=\frac53.
\]
取$c=3.0\times10^8$ m/s，
$v\Delta t'=120$ m，$v\Delta x'/c^2=-3.2\times10^{-7}$ s，
所以
\[
\boxed{\Delta x=0},\qquad
\boxed{\Delta t=3.0\times10^{-7}\,\mathrm s}.
\]
\crosscheck 两事件在$S$中同地发生，所以$3.0\times10^{-7}$ s是本题事件对的固有时间；$S'$中间隔$5\times10^{-7}$ s较长。

\sourceproblem{11.2.4}{290}{290}
\problemtext $\mu$子在自身静止系寿命$\tau_0=2.0\times10^{-6}$ s；
相对地球速率$0.988c$，求地球系测量的平均寿命。
\workings
由时间延缓
\[
\boxed{\tau=\gamma\tau_0
=\frac{2.0\times10^{-6}}{\sqrt{1-0.988^2}}
\approx1.295\times10^{-5}\,\mathrm s}.
\]
\crosscheck 运动粒子的地球系平均寿命应长于自身平均固有寿命。

\sourceproblem{11.2.5}{290}{290}
\problemtext 某不稳定粒子的固有平均寿命为实验室系测得寿命的$1/n$，
静质量$m_0$，求其实验室系动能。
\workings
$\tau_{\rm lab}=\gamma\tau_0$，故$\gamma=n$。
动能为
\[
\boxed{E_k=(\gamma-1)m_0c^2=(n-1)m_0c^2}.
\]
\crosscheck 当$n=1$对应粒子在该实验室系静止，动能为零。

\sourceproblem{11.2.6}{290}{290}
\problemtext 何种速度下粒子的相对论质量等于静质量的两倍？
何种速度下其动能恰等于静止能量？
\workings
两条件分别是$\gamma=2$与
$(\gamma-1)m_0c^2=m_0c^2$，
都要求$\gamma=2$，
\[
\boxed{\frac vc=\sqrt{1-\frac14}=\frac{\sqrt3}{2}}.
\]
所以两空相同。
\crosscheck 原题沿用教材的“运动质量”说法；现代惯用静质量$m_0$及相对论总能量$E=\gamma m_0c^2$，结论不变。
